\section{Physical Bridge Boundary}

The theorem established in this paper is finite relational
mathematics.

Native relational saturation is now an internal theorem of the
registered finite mechanics.

Quadratic contrast concentration is a second finite theorem-level
operation.

Physical interpretation belongs to a separate layer.

No result in the preceding sections is promoted to physical gravity,
black-hole physics, spacetime, or Planck-scale information theory
without an additional bridge theorem.

\subsection{Three layers}


The epistemic structure of Paper 14 is now
\[
\text{finite relational theorem}
\longrightarrow
\text{internal conditional extensions}
\longrightarrow
\text{physical bridge hypotheses}.
\]

The first layer contains native relational loading, registered-future
saturation, retained projection-relative distinction, and quadratic
contrast concentration.

The second contains proposed compositions such as clean release and
overcompression, together with stronger global transport
constructions.

The third requires an explicit mapping from finite theorem objects to
physical observables.

No physical noun promotes an internal definition.

Thus the phrases
\[
\text{relational saturation},
\qquad
\text{registration boundary},
\qquad
\text{clean release},
\qquad
\text{overcompression}
\]
remain internal terms unless a physical bridge is separately proved.

\subsection{Finite distinction support}

One possible physical extension begins from the hypothesis that an
independent physically realizable distinction-event requires finite
support.

Let
\[
r_E(\mathcal R)
\]
denote the number or rank of independent actionable distinction-events
realized in a physical region
\[
\mathcal R.
\]

A physical finite-density hypothesis might take the form
\[
r_E(\mathcal R)
\le
\rho_{E,\max}V_4(\mathcal R),
\]
where
\[
V_4(\mathcal R)
\]
is an independently defined physical four-volume.

A stronger dimensional conjecture might propose
\[
\rho_{E,\max}
\sim
\ell_P^{-4}.
\]

Equivalently, one might conjecture a minimum physical four-support of
order
\[
\ell_P^4
\]
per independent physically realizable distinction-event.

None of these statements is proved in Paper 14.

In particular, the paper does not assert
\[
1\text{ bit}=\ell_P^4.
\]

The relevant conjectural quantity is independent distinction density,
not arbitrary binary labeling.

\subsection{Event capacity is not static storage}

A four-volume density, if physically meaningful, cannot be identified
without argument with the static information capacity of a stationary
object.

A persistent physical system may occupy increasing four-volume as
duration accumulates while its instantaneous state capacity remains
unchanged.

Therefore any future
\[
\ell_P^4
\]
bridge must distinguish at least:

\begin{enumerate}[label=(\roman*)]

\item state capacity;

\item event or innovation capacity;

\item history capacity;

\item boundary-transmissible registration capacity.

\end{enumerate}

Paper 14 does not collapse these categories.

\subsection{Thalean gravity remains a bridge hypothesis}


The upgraded finite architecture separates two operations that earlier
versions left partially conflated:

\[
\text{native future selection}
\]
and
\[
\text{quadratic contrast concentration}.
\]

A future physical gravity hypothesis would therefore have to identify:

\begin{enumerate}[label=(\roman*)]

\item a physical state space or observable family;

\item a map from that physical structure to the finite relational
register;

\item a physical analogue of native future selection;

\item a physical analogue of contrast concentration, if required;

\item an explicit coupling law if both participate in one physical
process;

\item a derivation of the relevant aggregate and registered contrast;

\item and a demonstrated relation to gravitational observables.

\end{enumerate}

No such derivation is supplied here.

Therefore
\[
\text{finite relational saturation}
\not\Rightarrow
\text{physical gravity}.
\]

\subsection{No nonlinear graph deformation is introduced}

The finite mechanics used in Paper 14 is linear.

The graph and the quadratic operator are not deformed as a function
of loading.

The theorem does not propose a state-dependent metric
\[
G_{ij}(x),
\]
a nonlinear adjacency law, or a graph that physically bends under
source strength.

The exact operator remains
\[
Q=A^3+2A^2+2I.
\]

Any future physical gravity bridge must therefore respect the linear
finite substrate unless an independently justified extension is
introduced.

\subsection{Long range is not a separate selector}

Paper 14 also does not introduce a universal long-range gravitational
selector.

The internal QR architecture instead seeks long-range structure
through lawful accumulation over an admitted summation domain.

Schematically,
\[
\text{local admissible contributions}
\longrightarrow
\text{global compatibility}
\longrightarrow
\text{lawful sum}.
\]

A physical bridge would need to show how such a finite sum produces
the appropriate macroscopic gravitational behavior.

That bridge has not been derived.

\subsection{Black-hole interpretation}

Relational saturation may motivate a later black-hole hypothesis.

The internal mechanism would be:

\[
\text{finite complete distinction}
\]
together with
\[
\text{increasing aggregate dominance}
\]
and
\[
\text{decreasing outward individuation}.
\]

This provides an internal finite alternative to requiring unbounded
independent distinction density.

It does not prove that an astrophysical black hole realizes this
mechanism.

Accordingly, Paper 14 does not identify:

\begin{itemize}

\item contrast saturation with gravitational collapse;

\item a projection fibre with a black-hole interior;

\item a registration boundary with an event horizon;

\item common-mode dominance with mass density;

\item or finite relational degeneracy with a physical singularity.

\end{itemize}

These remain possible bridge hypotheses only.

\subsection{Horizon-like registration}

Within the internal mechanics one may define a horizon-like condition
relative to outward registration.

Let
\[
\Pi_{\mathrm{out}}
\]
denote a declared family of outward registrations.

A finite internal system is registration-saturated when additional
complete distinctions fail to produce additional independently
distinguishable outward registrations.

This can be stated schematically as
\[
r_{\mathrm{int}}
>
r_{\mathrm{out}}.
\]

That statement is internal.

A physical event horizon is a causal object in spacetime.

No theorem in Paper 14 identifies the two.

\subsection{No singularity theorem}

The finite contrast-concentration theorem contains no divergent
quantity.

Instead,
\[
\widehat Q^{\,k}
\]
remains finite for every
\[
k.
\]

Its centered component contracts, and its operator limit is the
finite rank-one projection
\[
P_0.
\]

This gives a finite algebraic alternative to divergence inside the
model.

It does not establish singularity avoidance in general relativity.

A physical singularity claim would require an independently derived
spacetime bridge and the appropriate causal, metric, and geodesic
analysis.

Therefore
\[
\text{finite saturation}
\not\Rightarrow
\text{physical singularity resolution}.
\]

\subsection{Black-hole entropy}

Paper 14 does not derive the Bekenstein-Hawking relation
\[
S_{\mathrm{BH}}
=
\frac{k_BA}{4\ell_P^2}.
\]

In particular, the conjectured four-support
\[
\ell_P^4
\]
for independent event distinction is not the same mathematical
quantity as an area-scaled black-hole entropy.

Any future relation between these ideas would have to explain:

\begin{enumerate}[label=(\roman*)]

\item why a four-dimensional event-capacity law produces an
area-scaled boundary law;

\item why the relevant coefficient is
\[
\frac14;
\]

\item which physical distinctions are being counted;

\item how the boundary transmits or registers them.

\end{enumerate}

None of those derivations is supplied here.

\subsection{No Hawking-radiation claim}

The clean-release language introduced in Section 11 is not identified
with Hawking radiation.

Clean release is an internal statement about coherent registered
output and surviving contrast.

Hawking radiation is a physical quantum-field-theoretic phenomenon
requiring a separate physical derivation.

No bridge between them is asserted.

\subsection{Cosmic microwave background}

A later QR interpretation may ask whether cosmological transparency
can be understood as an opening of a larger registration or
summation domain.

Paper 14 does not establish that interpretation.

The existing WMAP experiment tested a frozen native \(C_3\)
phase-structure detector and returned a formal NO-GO replication.

It therefore cannot be used here as evidence for a physical
saturation mechanism.

A future CMB saturation experiment would require a new frozen
observable and a new decision contract.

\subsection{Thalean time remains intrinsic finite time}

The established finite theory defines
\[
\tau_{\mathrm T}
=
\operatorname{wind}_{r_1}
\]
as the primitive additive oriented winding of the unique
\(G_{30}\)-visible registered-history cycle.

Paper 14 does not identify
\[
\tau_{\mathrm T}
\]
with:

\begin{itemize}

\item seconds;

\item proper time;

\item physical duration;

\item a Lorentzian temporal coordinate;

\item one quantum of action;

\item \(h\) or \(\hbar\);

\item or a spacetime interval.

\end{itemize}

No physical calibration is supplied by the saturation theorem.

\subsection{Physical falsifiability requirement}


A future physical bridge must make more than a qualitative analogy.

It must identify observables and predict a measurable relation that
could fail.

Because the finite theory now distinguishes native selection from
contrast concentration, the physical bridge must do the same.

A contrast-concentration prediction might have the form
\[
C_{\mathrm{phys}}(\lambda)
\quad
\text{remains large or increases}
\]
while
\[
\zeta_{\mathrm{phys}}(\lambda)
\quad
\text{decreases}.
\]

A native-selection prediction would require a separately defined
physical analogue of decreasing lawful future freedom.

A proposed physical cycle containing both must preregister the
coupling between them.

Without those mappings, the finite theorem remains mathematics rather
than physics.

\subsection{Explicit nonclaims}

Paper 14 does not establish:

\begin{enumerate}[label=(\roman*)]

\item physical information density;

\item one bit per Planck four-volume;

\item physical gravity;

\item Newtonian gravity;

\item general relativity;

\item the equivalence principle;

\item physical spacetime;

\item black-hole interiors;

\item physical event horizons;

\item singularity avoidance;

\item black-hole entropy;

\item holography;

\item Hawking radiation;

\item a physical CMB saturation signal;

\item or a Standard Model embedding.

\end{enumerate}

\subsection{Bridge keeper}


The relationship among the layers can be summarized as follows:

\begin{quote}
Finite relational mathematics can support internal mechanical
extensions.  Internal mechanics can motivate physical hypotheses.
Physical claims do not follow without an explicit physical bridge.
\end{quote}

Paper 14 now establishes both a native finite relational-saturation
theorem and a separate quadratic contrast-concentration theorem.

Clean-release composition, stronger global inter-face transport, and
all physical interpretations remain additional hypotheses.

Whether nature realizes any analogue of this finite architecture
remains an open physical question.
